% ══════════════════════════════════════════════════════════════
%  SECTION 4: EMPIRICAL STRATEGY
% ══════════════════════════════════════════════════════════════

\section{Empirical Strategy}
\label{sec:strategy}

\subsection{The Identification Challenge}

The directive did not produce a clean experiment. FIFA's announcement
was global, the IFAB endorsement was universal, and the World Cup
demonstration was visible to every federation, referee assessor, and
television audience in Europe simultaneously. What followed was not
selective adoption by five leagues but a near-universal norm shift in
which 13 of 14~leagues with stoppage-time data increased their
enforcement. The binary treatment--control distinction we impose---Big
Five top-flight leagues as treated, all others as controls---is a
pragmatic simplification of a continuous phenomenon.

Partial treatment of the control group is the central identification
challenge. Several control leagues increased stoppage time
substantially: the Championship added 1.03~minutes ($d = 0.59$),
2.~Bundesliga added 1.04~minutes ($d = 0.59$), and the Scottish
Premiership added 1.47~minutes ($d = 0.83$). These leagues share
institutional infrastructure with their treated counterparts---referee
pools overlap between divisions, assessors circulate the same
directives, and media coverage of the World Cup's extended matches
reached every football culture in Europe. A naive comparison of
treated and control groups therefore understates the true effect of
the directive, because the control group absorbed a large share of
the same treatment.

Two interpretations of the control-group increase compete. First, the
directive may have leaked through institutional channels: shared
referee development programs, cross-divisional assessor networks, and
ubiquitous media exposure transmitted the new enforcement norm to
leagues that were not formal recipients. Under this interpretation,
the control-group shift is itself evidence of the directive's reach,
and the binary DiD captures only the differential between the top and
middle of a dose gradient---a conservative lower bound on the full
effect. Second, there may have been a secular trend in stoppage time
unrelated to the directive, driven by the proliferation of VAR
reviews, changes in match tempo, or evolving attitudes toward time
management. Under this interpretation, the control-group shift
reflects a pre-existing trajectory that would have occurred without
FIFA's intervention, and the DiD estimate is contaminated by a
violation of parallel trends.

Our strategy addresses this ambiguity through multiple designs rather
than a single dispositive specification. The pooled binary DiD
provides a lower bound under the first interpretation. Within-country
pairs absorb country-level secular trends, testing whether the
residual treated--control gap survives after removing shared national
shocks. The dose-response framework reframes the analysis entirely,
treating the binary DiD as one point on a gradient that is itself
the object of interest.

\subsection{Primary: Binary Difference-in-Differences}

The primary specification pools all 15~leagues with stoppage-time data
across 12~seasons. We estimate:
%
\begin{equation}\label{eq:did}
  Y_{ilt} = \alpha + \beta \cdot (\text{Post}_t \times \text{Treated}_l)
  + \mu_l + \gamma_t + \varepsilon_{ilt}
\end{equation}
%
where $Y_{ilt}$ is the outcome for match $i$ in league $l$ and season
$t$; $\text{Post}_t$ equals one for seasons beginning in 2023--24 or
later; $\text{Treated}_l$ equals one for the five Big~Five top-flight
leagues; $\mu_l$ and $\gamma_t$ are league and season fixed effects.
The coefficient $\beta$ captures the average change in the treated--control
gap after the directive. Because adoption was effectively simultaneous
across treated leagues---domestic seasons began between August 11 and
August 19, 2023---this is a standard two-period DiD with a common
treatment date, avoiding the staggered-adoption complications that
have motivated recent methodological work
\citep{goodmanbacon2021difference, callaway2021difference,
sunAbraham2021, borusyak2024revisiting}. We nonetheless verify that
our estimates are robust to all three heterogeneity-robust estimators
as a specification check.

Inference requires care. With 15~league-level clusters and treatment
assigned at the league level, conventional cluster-robust standard
errors may over-reject. We report wild cluster bootstrap $p$-values
with 999~Rademacher weight replications throughout
\citep{cameron2008bootstrap}. Following \citet{webb2023reworking}, we
also report the six-point Webb weights as a further robustness check,
since Rademacher weights can behave poorly with very few treated
clusters. In practice, the two weighting schemes produce similar
$p$-values for our main specifications.

The coefficient $\beta$ should be interpreted as a conservative lower
bound on the directive's effect. Control-group contamination biases
the estimate toward zero: if control leagues also increased stoppage
time in response to the directive, the treated--control gap shrinks
relative to the counterfactual in which control leagues were truly
untreated. The magnitude of this attenuation is substantial---control
leagues increased stoppage time by an average of 1.30~minutes,
compared to 1.57~minutes in treated leagues---implying that the
binary DiD captures roughly the upper quartile of the enforcement
distribution minus the median, not the full treatment effect.

\subsection{Co-Primary: Within-Country Pairs}

To tighten identification, we restrict the sample to countries that
contribute both a treated top-flight league and an untreated second
division: England (E0/E1), Germany (D1/D2), France (F1/F2), and Italy
(I1/I2). Spain is excluded because Segunda~Divisi\'{o}n data are
unavailable after 2022--23. This yields eight league-level clusters
across four countries.

The within-country specification is:
%
\begin{equation}\label{eq:within}
  Y_{ilct} = \alpha + \beta \cdot (\text{Post}_t \times \text{TopFlight}_l)
  + \gamma_{ct} + \mu_l + \varepsilon_{ilct}
\end{equation}
%
where $\gamma_{ct}$ denotes country$\times$season fixed effects and
$\mu_l$ denotes league fixed effects. The country$\times$season terms
absorb any shock common to both divisions within a country in a given
season---including VAR adoption timelines, national media cycles,
federation-level policy changes, and macroeconomic conditions that
might independently affect match play. What remains is the
within-country, within-season difference between the top-flight and
second-division response.

This design substantially tightens identification but reduces
precision. Eight clusters afford limited degrees of freedom, and the
wild cluster bootstrap with so few clusters approaches the boundary
of its asymptotic validity. We report both standard cluster-robust
and wild bootstrap $p$-values, noting that the latter should be
interpreted cautiously.

The Italy pair warrants specific discussion. Serie~A shows a muted
response to the directive ($+0.66$~minutes, $d = 0.40$), while
Serie~B shows essentially no response ($-0.05$~minutes,
$d = -0.04$). The bilateral within-country estimate for Italy is
therefore negative: the second division, supposedly the cleaner
control, did not shift, but nor did the top-flight shift much. In
the \citet{goodmanbacon2021difference} decomposition, the Italy pair
pulls the within-country estimate downward. The within-country design
therefore relies effectively on three country pairs---England,
Germany, and France---where the top-flight response meaningfully
exceeded the second-division response. We report the four-pair
estimate as the primary within-country result but note that
excluding Italy strengthens the coefficient and tightens confidence
intervals.

\subsection{The Dose-Response Interpretation}

The binary treated--control distinction obscures the most striking
feature of the data: enforcement intensity varies continuously across
leagues, forming a gradient that is itself informative about the
directive's transmission mechanism. Table~\ref{tab:enforcement}
documents this gradient. At the top, the Premier League ($d = 1.13$),
Ligue~1 ($d = 1.05$), and the Bundesliga ($d = 0.92$) shifted by
approximately a full standard deviation. In the middle, the Scottish
Premiership ($d = 0.83$), the Championship
($d = 0.59$), and the 2.~Bundesliga ($d = 0.59$) shifted nearly as
much as some treated leagues. At the bottom, Serie~A ($d = 0.40$)
and the Eredivisie ($d = 0.23$) responded modestly, and Serie~B ($d = -0.04$) did not
respond at all.

What predicts a league's position on this gradient? The strongest
correlates are institutional: Big~Five status (which implies direct
representation on UEFA and FIFA committees), the presence of a
professional refereeing structure with full-time officials, the
intensity of domestic media coverage of the World Cup's extended
matches, and---critically---whether the league shares a referee
development pipeline with a Big~Five top-flight division. The
Championship and 2.~Bundesliga, whose referees are trained alongside
Premier League and Bundesliga officials, shifted almost identically
to their top-flight counterparts. The Scottish Premiership, whose
referees operate independently but whose media market is saturated
by English football coverage, shifted even more than some formally
treated leagues.

We do not estimate a formal dose-response regression. The dose is
measured at the league level, leaving only 14~observations (one
league lacks post-directive stoppage data); parametric estimation
with so few clusters would be unreliable regardless of the functional
form assumed. Instead, we present the dose gradient descriptively
and use it to interpret the binary DiD: $\hat{\beta}$ captures the
difference between the top of the gradient (the Big~Five minus
Serie~A) and the middle (the contaminated control group), not the
full distance from the pre-directive equilibrium. Readers who want
a point estimate of the directive's total effect should look to the
treated-group pre--post change (1.57~minutes) rather than the DiD
coefficient (0.74~minutes), recognizing that the former lacks a
counterfactual and the latter is attenuated.

\subsection{Pre-Trends and Their Limits}

Honest assessment of the parallel-trends assumption requires
acknowledging what the data show. The event study rejects the null
of parallel pre-trends: a joint $F$-test of pre-treatment
season-by-treatment interactions yields $F = 9.007$ ($p < 0.001$).
Treated and control leagues were on diverging stoppage-time
trajectories before the directive. We do not minimize this finding.

To formalize what the pre-trend violation implies for post-treatment
inference, we apply the \citet{rambachan2023more} HonestDiD
framework. The breakdown value is $\bar{M}^* = 0.11$: if one
permits the post-treatment differential trend to deviate from a
linear extrapolation of the pre-treatment trend by more than
0.11~minutes per period, the confidence interval for $\beta$
includes zero. The individual post-period breakdown values are
negative ($-0.31$, $-0.67$, and $-0.51$ for 2023--24, 2024--25,
and 2025--26 respectively), meaning that under \textit{any}
positive extrapolation of pre-existing trends, the smoothed
confidence intervals include zero. By the strict criterion of
\citet{roth2022pretest}---requiring robustness to trend
extrapolation---the binary DiD does not survive.

We proceed for four reasons, none of which requires dismissing the
pre-trend evidence.

First, the within-country specification absorbs country-level secular
trends. The country$\times$season fixed effects in
Equation~\ref{eq:within} remove any trajectory shared by both
divisions within a country, including differential VAR adoption
timelines and country-specific changes in match tempo. The
within-country estimate ($\hat{\beta} = 0.87$, $p = 0.033$)
survives, and its pre-treatment coefficients are closer to zero than
those in the pooled specification, suggesting that much of the
pre-trend failure reflects cross-country heterogeneity rather than
a violation of the identifying assumption within country pairs.

Second, the mechanism is directly observable. Unlike most
difference-in-differences applications, where the treatment is
inferred from policy variation and the counterfactual is
unobservable, the stoppage-time directive was publicly announced,
its implementation is recorded in every match report, and referee
compliance can be measured individual by individual. The question is
not whether the directive caused a change in referee behavior---it
manifestly did---but whether the DiD correctly measures the
magnitude. The pre-trend violation speaks to magnitude, not to the
existence of an effect.

Third, the pre-trend failure has a plausible structural explanation.
Cross-league trajectories in stoppage time diverged in the
pre-directive period for identifiable reasons: VAR was adopted on
different timelines across leagues, match tempo evolved
heterogeneously, and culturally specific attitudes toward time
management (the Italian tradition of \textit{catenaccio}-era
time-wasting, the English norm of minimal stoppage) generated
cross-country drift. These sources of divergence are real but do
not imply that treated leagues would have continued to diverge from
controls absent the directive. The within-country design, which
removes national-level trends, is precisely the specification that
addresses this concern, and it is the specification that survives.

Fourth, the displacement result---total goals unchanged, with
temporal redistribution toward stoppage time---does not depend on
the cross-league comparison at all. It holds within the treated
leagues alone, comparing the pre- and post-directive distribution
of goals across match minutes. No control group is needed, and no
parallel-trends assumption is invoked.

How, then, should readers interpret the results? Those who require
$\bar{M}^* > 0$ for the pooled DiD to be credible should treat
our estimates as what \citet{roth2022pretest} call ``precise
descriptive accounting of what changed when enforcement norms
shifted''---a characterization we embrace rather than resist. The
paper's contribution does not rest on any single specification being
dispositive. It rests on the conjunction of a directly observable
mechanism, a dose gradient consistent with institutional
transmission, a within-country estimate that survives tighter
identification, and a displacement result that is robust to the
specification choices that affect the binary DiD. The stoppage-time
directive changed European football. The empirical question is how
precisely we can measure by how much, and the answer depends on how
much weight one places on the parallel-trends assumption in a
setting where partial treatment of the control group was inevitable
\citep{rothetal2023, imbens2009recent}.
